\section{Methodology}

Esta sección describe el diseño experimental preliminar. Cada decisión busca que los
resultados sean reproducibles y libres de sesgos por el uso indebido de datos de prueba.

\subsection{Base de Datos}

Se utiliza la base de datos AT\&T Database of Faces \cite{samaria1994parameterisation}, la cual
contiene 400 imágenes de 40 personas en escala de grises; es decir, 10 imágenes por persona.
Cada imagen tiene $92 \times 112$ píxeles (columnas $\times$ filas) con 256 niveles de gris
(del 0 al 255). Las personas aparecen en posición frontal sobre un fondo homogéneo y, para
algunas de ellas, las imágenes se tomaron en momentos distintos, con variaciones de
iluminación, expresión facial y detalles faciales.

Al representar cada imagen como una matriz $A \in \mathbb{R}^{112 \times 92}$, se tiene
$mn = 10304$; por tanto, la matriz de covarianza $C$ tendría más de $10^8$ componentes,
lo cual hace inviable su cálculo directo y motiva el uso de la SVD.

Se asume que cada imagen contiene un rostro recortado y alineado. Las intensidades se leen
como enteros de 0 a 255 y se convierten a punto flotante de doble precisión antes de
centrar; no se normalizan a $[0,1]$, pues escalar todos los valores por igual no altera qué
imagen queda más cercana.


\subsection{Protocolo Experimental}

Para cada persona se eligen al azar $n_{\text{ent}}$ imágenes de entrenamiento y las
$10 - n_{\text{ent}}$ restantes de prueba, de modo que $N = 40\,n_{\text{ent}}$. Ninguna
imagen de prueba participa en el cálculo de $\bar{\mathbf{x}}$, de la SVD ni de las
coordenadas de entrenamiento.
\begin{itemize}
    \item \textbf{División principal:} $n_{\text{ent}} = 5$; es decir, $N = 200$ imágenes
    de entrenamiento y 200 de prueba. Es una división balanceada, con un conjunto de prueba
    del mismo tamaño que el de entrenamiento.
    \item \textbf{Repeticiones:} $R = 20$ divisiones aleatorias con semillas fijas.
    Se reporta la media y la desviación estándar, pues con 200 imágenes de prueba una
    sola división puede favorecer o perjudicar el resultado por azar.
    \item \textbf{Tamaño del entrenamiento:} $n_{\text{ent}} \in \{2, 3, \dots, 8\}$.
    Esto permite contrastar la afirmación de que Eigenfaces opera con pocos datos.
\end{itemize}

\textbf{Elección de $k$.} Para no usar las imágenes de prueba al elegir $k$, se reporta la
curva completa de la tasa de reconocimiento para $k = 1, \dots, N-1$. Además, el valor de
operación se fija con el criterio de varianza explicada: el menor $k$ tal que
$\text{VE}(k) \geq \tau$, con $\tau \in \{0.90, 0.95, 0.99\}$.
Esta regla depende solo de los datos de entrenamiento. Si las coordenadas se calculan con
las $N-1$ eigenfaces, la distancia que usa solo las primeras $k$ componentes es
\[
d_j^2(k) = \sum_{i=1}^{k}\big([\mathbf{y}]_i - [\mathbf{y}_j]_i\big)^2
\]
por lo que la curva completa se obtiene con una sola SVD por división.


\subsection{Métricas y Líneas Base}

Se define la \textbf{tasa de reconocimiento} como:
\[
\text{TR} = \frac{\text{imágenes de prueba identificadas correctamente}}
{\text{total de imágenes de prueba}}
\]
Se reportan también $\text{VE}(k)$ y el error relativo de reconstrucción
$\|\tilde{X} - \tilde{X}_k\|_F / \|\tilde{X}\|_F$ en función de $k$, que conectan los
experimentos con la proposición sobre el error de aproximación y el teorema de
Eckart--Young. Para interpretar $\text{TR}$ se usan dos líneas base:
\begin{itemize}
    \item \textbf{Azar:} con 40 personas, $\text{TR} = 1/40 = 2.5\,\%$.
    \item \textbf{Vecino más cercano en píxeles:} el mismo criterio de distancia mínima,
    pero en $\mathbb{R}^{mn}$, sin reducción. Permite medir qué se gana o se pierde al
    pasar de $mn$ a $k$ dimensiones.
\end{itemize}

Se reporta también el tiempo de cómputo del ajuste (cálculo de $\bar{\mathbf{x}}$ y de
la SVD) y del reconocimiento de una imagen, medido en el mismo equipo y con las mismas
dependencias, lo que sustenta la afirmación sobre el costo computacional.


\subsection{Implementación y Verificación}

El sistema se implementa en Python con NumPy; Matplotlib se usa solo para las figuras. Se
emplea la SVD delgada de $\tilde{X}$ (\texttt{numpy.linalg.svd} con
\texttt{full\_matrices=False}): como $N \ll mn$, evita construir $C$ y la matriz $V$
completa, y basta con sus primeras $N$ columnas, pues $\tilde{X}\mathbf{v}_i = \mathbf{0}$
para $i > r$. El código se organiza en cuatro módulos: carga y partición de los datos,
modelo (ajuste, proyección y clasificación), evaluación de los experimentos y pruebas.
Las dependencias se fijan por versión y las semillas son explícitas.

Para unir la teoría con el código, se incluyen pruebas automáticas que verifican con
tolerancia numérica los resultados del marco teórico:
\begin{itemize}
    \item $V^TV = I$ y $\tilde{X}V = U\Sigma$;
    \item $\|\tilde{X} - \tilde{X}_k\|_F^2 = \sum_{i=k+1}^{r}\sigma_i^2$;
    \item $\|\mathbf{y} - \mathbf{y}'\| \leq \|\mathbf{x} - \mathbf{x}'\|$;
    \item $\mathbf{1}^T\tilde{X} = \mathbf{0}$, es decir, $r \leq N-1$;
    \item $\sigma_i^2$ coincide con los eigenvalores de $\tilde{X}^T\tilde{X}$ en un
    subconjunto pequeño de píxeles, donde sí es factible calcularlos.
\end{itemize}
Como contraste adicional, las proyecciones se comparan, salvo el signo de cada eigenface,
con las de \texttt{scikit-learn} únicamente dentro de las pruebas; esta biblioteca no forma
parte del sistema.